% GENERATED FILE. Edit canonical lessons, then run npm run generate:books.
% canonical-source-hash: fnv1a64:c92d46ed6df437ec
% canonical-lessons: LA-C224-lacerta, LA-C224-aper, LA-C224-caper, LA-C224-agnus, LA-C224-vitulus

\chapter{Lizard, Wild boar, Billy goat, Lamb, Calf}
\label{ch:la-lizard-wild-boar-billy-goat-lamb-calf}
\hlchaptermodality{224}

\begin{chapteropening}
\textbf{By the end of this chapter:} \emph{I can say a lizard, a wild boar, a billy goat, a lamb and a calf.}

Complete the last lesson of chapter 224: I can say a lizard, a wild boar, a billy goat, a lamb and a calf.
\end{chapteropening}

\section[lacerta, lacertae]{lacerta, lacertae — a lizard}

\label{lesson:LA-C224-lacerta}



\begin{quote}
Before the new one: say the Latin for a thorn, then the Latin for a crop, a cornfield.
\end{quote}

\subsection*{You'll want to know: lacerta, lacertae}
\textbf{lacerta, lacertae} — \textquotedblleft{}a lizard\textquotedblright{}.

\begin{cousinweb}
\textbf{lacerta} has no sure origin; it may be linked with \emph{lacertus}, \textquotedblleft{}the upper arm, a muscle\textquotedblright{}, but that is uncertain. English \emph{lizard} comes from it through French.
\end{cousinweb}

\begin{grammarlens}[title={What you've built}]
One more word for this chapter.
\end{grammarlens}

\subsection*{Your turn}
\begin{itemize}
\raggedright
  \item \emph{Say it:} \emph{lacerta, lacertae}
  \item \emph{Say it:} \emph{lacerta, lacertae}, once more
  \item \emph{Say it:} say \emph{lacerta, lacertae}
  \item \emph{Recall:} say the Latin for a thorn, then the Latin for a crop, a cornfield, then say \emph{lacerta, lacertae} again
\end{itemize}

\begin{tcolorbox}[breakable,colback=teal!4,colframe=teal!35!black,title={Before you move on}]
What does lacerta, lacertae mean? (\textquotedblleft{}A lizard\textquotedblright{}.) Say it once more.
\end{tcolorbox}
\section[aper, aprī]{aper, aprī — a wild boar}

\label{lesson:LA-C224-aper}



\begin{quote}
Before the new one: say the Latin for a crop, a cornfield, then the Latin for a lizard.
\end{quote}

\subsection*{You'll want to know: aper, aprī}
\textbf{aper, aprī} — \textquotedblleft{}a wild boar\textquotedblright{}.

\begin{cousinweb}
\textbf{aper} is inherited, usually traced to Proto-Indo-European \emph{*h1epr-}, \textquotedblleft{}boar\textquotedblright{} (EtymDB 2.1); German \emph{Eber}, \textquotedblleft{}boar\textquotedblright{}, is a cousin (Lewis \& Short).
\end{cousinweb}

\begin{grammarlens}[title={What you've built}]
One more word for this chapter.
\end{grammarlens}

\subsection*{Your turn}
\begin{itemize}
\raggedright
  \item \emph{Say it:} \emph{aper, aprī}
  \item \emph{Say it:} \emph{aper, aprī}, once more
  \item \emph{Say it:} say \emph{aper, aprī}
  \item \emph{Recall:} say the Latin for a crop, a cornfield, then the Latin for a lizard, then say \emph{aper, aprī} again
\end{itemize}

\begin{tcolorbox}[breakable,colback=teal!4,colframe=teal!35!black,title={Before you move on}]
What does aper, aprī mean? (\textquotedblleft{}A wild boar\textquotedblright{}.) Say it once more.
\end{tcolorbox}
\section[caper, caprī]{caper, caprī — a billy goat}

\label{lesson:LA-C224-caper}



\begin{quote}
Before the new one: say the Latin for a lizard, then the Latin for a wild boar.
\end{quote}

\subsection*{You'll want to know: caper, caprī}
\textbf{caper, caprī} — \textquotedblleft{}a billy goat\textquotedblright{}.

\begin{cousinweb}
\textbf{caper} is inherited, usually traced to Proto-Indo-European \emph{*kapros}, \textquotedblleft{}buck, he-goat\textquotedblright{} (EtymDB 2.1); Greek \emph{kapros}, \textquotedblleft{}wild boar\textquotedblright{}, is a cousin (Lewis \& Short).
\end{cousinweb}

\begin{grammarlens}[title={What you've built}]
One more word for this chapter.
\end{grammarlens}

\subsection*{Your turn}
\begin{itemize}
\raggedright
  \item \emph{Say it:} \emph{caper, caprī}
  \item \emph{Say it:} \emph{caper, caprī}, once more
  \item \emph{Say it:} say \emph{caper, caprī}
  \item \emph{Recall:} say the Latin for a lizard, then the Latin for a wild boar, then say \emph{caper, caprī} again
\end{itemize}

\begin{tcolorbox}[breakable,colback=teal!4,colframe=teal!35!black,title={Before you move on}]
What does caper, caprī mean? (\textquotedblleft{}A billy goat\textquotedblright{}.) Say it once more.
\end{tcolorbox}
\section[agnus, agnī]{agnus, agnī — a lamb}

\label{lesson:LA-C224-agnus}



\begin{quote}
Before the new one: say the Latin for a wild boar, then the Latin for a billy goat.
\end{quote}

\subsection*{You'll want to know: agnus, agnī}
\textbf{agnus, agnī} — \textquotedblleft{}a lamb\textquotedblright{}.

\begin{cousinweb}
\textbf{agnus} is inherited, usually traced to Proto-Indo-European \emph{*h2egwnos}, \textquotedblleft{}lamb\textquotedblright{} (EtymDB 2.1); Greek \emph{amnos}, \textquotedblleft{}lamb\textquotedblright{}, is a cousin (Lewis \& Short), and so is English \emph{yean}, \textquotedblleft{}to give birth to lambs\textquotedblright{}.
\end{cousinweb}

\begin{grammarlens}[title={What you've built}]
One more word for this chapter.
\end{grammarlens}

\subsection*{Your turn}
\begin{itemize}
\raggedright
  \item \emph{Say it:} \emph{agnus, agnī}
  \item \emph{Say it:} \emph{agnus, agnī}, once more
  \item \emph{Say it:} say \emph{agnus, agnī}
  \item \emph{Recall:} say the Latin for a wild boar, then the Latin for a billy goat, then say \emph{agnus, agnī} again
\end{itemize}

\begin{tcolorbox}[breakable,colback=teal!4,colframe=teal!35!black,title={Before you move on}]
What does agnus, agnī mean? (\textquotedblleft{}A lamb\textquotedblright{}.) Say it once more.
\end{tcolorbox}
\section[vitulus, vitulī]{vitulus, vitulī — a calf}

\label{lesson:LA-C224-vitulus}



\begin{quote}
Before the new one: say the Latin for a billy goat, then the Latin for a lamb.
\end{quote}

\subsection*{You'll want to know: vitulus, vitulī}
\textbf{vitulus, vitulī} — \textquotedblleft{}a calf\textquotedblright{}.

\begin{cousinweb}
\textbf{vitulus} is usually traced to Proto-Indo-European \emph{*wetelo-}, \textquotedblleft{}yearling\textquotedblright{}, from \emph{*wet-}, \textquotedblleft{}year\textquotedblright{} (EtymDB 2.1): a calf as a one-year-old. The name \emph{Italia}, \textquotedblleft{}Italy\textquotedblright{}, is often explained as \textquotedblleft{}land of calves\textquotedblright{} (Lewis \& Short).
\end{cousinweb}

\begin{grammarlens}[title={What you've built}]
One more word for this chapter.
\end{grammarlens}

\subsection*{Your turn}
\begin{itemize}
\raggedright
  \item \emph{Say it:} \emph{vitulus, vitulī}
  \item \emph{Say it:} \emph{vitulus, vitulī}, once more
  \item \emph{Say it:} say \emph{vitulus, vitulī}
  \item \emph{Recall:} say the Latin for a billy goat, then the Latin for a lamb, then say \emph{vitulus, vitulī} again
\end{itemize}

\begin{tcolorbox}[breakable,colback=teal!4,colframe=teal!35!black,title={Before you move on}]
What does vitulus, vitulī mean? (\textquotedblleft{}A calf\textquotedblright{}.) Say it once more.
\end{tcolorbox}
